% Cabeçalho LaTeX comum aos livros (inserido por pandoc -H).
\usepackage{microtype}      % justificação mais fina
\usepackage{fancyhdr}
\pagestyle{fancy}
\fancyhf{}
\fancyfoot[C]{\small\thepage}
\renewcommand{\headrulewidth}{0pt}
\fancypagestyle{plain}{\fancyhf{}\fancyfoot[C]{\small\thepage}\renewcommand{\headrulewidth}{0pt}}
\raggedbottom
\setlength{\parindent}{1.2em}
\setlength{\parskip}{0pt}

% Títulos na própria fonte do texto (sem sans-serif), centralizados
\usepackage{titlesec}
\titleformat{\chapter}[display]{\normalfont\centering\LARGE\bfseries}{}{0pt}{}
\titlespacing*{\chapter}{0pt}{20pt}{30pt}
\titleformat{\section}{\normalfont\Large\bfseries\centering}{}{0pt}{}
\titlespacing*{\section}{0pt}{2.2ex plus .5ex}{1.6ex}
\titleformat{\subsection}{\normalfont\large\bfseries\centering}{}{0pt}{}
\titlespacing*{\subsection}{0pt}{1.8ex plus .4ex}{1.2ex}
\titleformat{\subsubsection}{\normalfont\normalsize\bfseries\itshape\centering}{}{0pt}{}

% Folha de rosto tipográfica (\fontedorosto: fonte dos títulos do livro, quando ele tem estilo próprio)
\providecommand{\fontedorosto}{}
\makeatletter
\renewcommand{\maketitle}{%
  \begin{titlepage}\centering
    \vspace*{0.22\textheight}
    {\fontedorosto\fontsize{26}{32}\selectfont\bfseries\@title\par}
    \vspace{3em}
    {\Large\@author\par}
    \vfill
  \end{titlepage}}
\makeatother
